\section{Galois categories}
\label{V.5}
% original p. 127

\begin{definition}
\label{V.5.1}
A Galois category is
a category $\cal{C}$ equivalent to a category
$\cal{C}(\pi)$, where $\pi$ is a compact group, projective limit
of finite groups (\ie totally disconnected).

For the definition of~$\cal{C}(\pi)$, \cf the beginning of \No \Ref{V.4}. By
Theorem~\Ref{V.4.1}, $\cal{C}$ is Galois if and only if
it satisfies conditions (G 1) to (G 3), and if there exists a
functor $F$ from
$\cal{C}$
into the category of finite sets satisfying conditions
(G~4) to (G~6) (\ie which is \emph{exact} and \emph{conservative},
in a general terminology). Such a functor will be
called a \emph{fundamental functor}
of the Galois category $\cal{C}$\kern1pt\footnote{It seems preferable to adopt the more suggestive term
``fiber functor''.};
it is pro-representable by a pro-object which we shall denote by~$P_F$;
a pro-object $P$ such that the associated functor $F$ is fundamental
is called a \emph{fundamental pro-object}.
In this way, the category of fundamental functors on~$\cal{C}$ is anti-equivalent to the category of
fundamental pro-objects; if $F$ and $P$ correspond to each other, the group
$\Aut F$ is therefore isomorphic to the opposite of the group $\Aut P$,
hence the group denoted $\pi$ in the preceding no.\ is
none other than $\Aut P$. Let us recall that in the preceding no.\ we
constructed, starting from a \emph{given} fundamental functor
$F$, an equivalence of~$\cal{C}$ with
$\cal{C}(\pi)$ (where $\pi=\Aut(F)$) which transforms the functor $F$
into the canonical functor from~$\cal{C}(\pi)_x$ into the category of
finite sets. In this typical case $\cal{C}=\cal{C}(\pi)$, $F=$
canonical functor, the fundamental pro-object associated with $F$
is none other than the projective system of the discrete quotients
$\pi_i$ of~$\pi$.
\end{definition}

It may be useful to make explicit the category of pro-objects of
$\cal{C}(\pi)$. One finds:

\begin{proposition}
\label{V.5.2}
The category $\Prodash\cal{C}(\pi)$
is canonically equivalent to the category
$\cal{C}'(\pi)$ of spaces, with topological group $\pi$
of operators, which are compact and totally disconnected.
\end{proposition}

Since the latter contains $\cal{C}(\pi)$ as a
full subcategory (corresponding to the compact discrete spaces with
operators) and since projective limits
exist in it, one has in any case a canonical functor $g\colon
\Prodash\cal{C}(\pi) \to \cal{C}'(\pi)$, associating with the projective system
$Q=(Q_i)$
the object $X=\varprojlim_{i} Q_i$ of~$\cal{C}'(\pi)$. To define
a functor in the opposite direction, it suffices to define a contravariant
functor from~$\cal{C}'(\pi)$ into the category of left exact functors
$\cal{C}\to \Ens$, and we shall take for every
$X\in\cal{C}'(\pi)$ the functor $h(X)(E)= \Hom(X, E)$
% original p. 128
(the $\Hom$ being taken in $\cal{C}'(\pi)$). It is immediate
by definition that the functors $h$ and $g$ are adjoint to each
other, and that $hg$ is canonically isomorphic to the identity functor
of~$\Prodash\cal{C}(\pi)$. It remains to prove (in order to establish
that $g$ and $h$ are quasi-inverse to each other) that every object of
$\cal{C}'(\pi)$ is isomorphic to an object of the form $g(Q)$,
where $Q\in\Prodash{\cal C}(\pi)$, in other words: \emph{every
space $X$ with topological group $\pi$ of operators, which is
compact and totally disconnected, is isomorphic to a projective
limit of finite discrete spaces with operators}. Since $X$
is the projective limit of its finite discrete quotients (as a
topological space without operators), we are reduced to
showing that in the set of these quotients, there is a
cofinal system which is invariant under $\pi$. For this it suffices to show
that for such a quotient $X'$, the set of transforms of this
quotient by the operations of~$\pi$ is finite, (one will then take
the sup of the said transforms, which will be an invariant quotient
majorizing $X'$). Or again, that there is an open invariant subgroup
$\pi'$ of~$\pi$ such that the elements of this subgroup
leave $X'$ invariant. Now $X'$ corresponds to a finite partition of~$X$ into
open sets $X_i$. By reason of continuity and of
compactness of~$\pi$, there exists a neighborhood $V$ of
the neutral element of~$\pi$ such that $s\in V$ implies
$s\cdot X_i\subset X_i$ for every~$i$, hence $s$ leaves~$X'$ invariant. Now we know
that the open invariant subgroups of~$\pi$ form a
fundamental system of neighborhoods of the neutral element, which
completes the proof.

Let us note that one sees even more simply that the category
$\Ind \cal{C}(\pi)$ is canonically equivalent to the
category of sets on which $\pi$ operates
continuously. We shall not need this here.

\begin{proposition}
\label{V.5.3}
Let $\cal{C}$ be a Galois category, $F$ a fundamental
functor
on~$\cal{C}$,
$P=(P_i)$ the associated pro-object, normalized in the usual
way. Let
$X\in\cal{C}$;
for $X$ to be connected, it is necessary and sufficient that $\pi$ operate
transitively on~$E=F(X)$.
\end{proposition}

We are reduced to the typical case $\cal{C}=\cal{C}(\pi)$, $F=$ canonical
functor, where it is trivial.

\begin{corollary}
\label{V.5.4}
Equivalent conditions on~$X$: \textup{(i)} $X$ is connected and
$\not\simeq\emptyset_{\cal{C}}$ \textup{(ii)} the group $\pi$ is
transitive on~$E=F(X)$, and $F(X)\ne\emptyset$ \textup{(iii)} $X$ is
isomorphic to a $P_i$.
\end{corollary}

The equivalence
% original p. 129
of (i) and (iii) also already follows easily from
\No \Ref{V.4},~e).

\begin{proposition}
\label{V.5.5}
Let $Q=(Q_i)_{i\in I}$ be a pro-object
of~$\cal{C}$,
normalized in the usual way, and let $G$ be the
corresponding functor $G(X)=\Hom(Q, X)$ from~$\cal{C}$
(into~$\Ens$).
The following conditions are equivalent:
\begin{enumerate}
\item[(i)] $G$ commutes with finite direct sums
\item[(ii)] $G$ commutes with the sum of two objects
\item[(iii)] The $Q_i$ are connected and
$\not\simeq\emptyset_\cal{C}$
\item[(iv)] $Q$ is isomorphic to $\pi/H$, where $H$ is a
closed subgroup of~$\pi$.
\item[(v)] The functor $G$ is isomorphic to the functor $E\mto
E^H$ (set of invariants under $H$) defined by a closed
subgroup $H$ of~$\pi$.
\end{enumerate}
\end{proposition}

N.B. in the statement of (iv) and (v), one assumes chosen a
fundamental functor, making it possible to identify $\cal{C}$ with the
category~$\cal{C}(\pi)$.

\subsubsection*{Proof} We may assume
$\cal{C}=\cal{C}(\pi)$. The implication (i)$\To$(ii) is
trivial, (ii)$\To$(iii) is proved like
property d) of \No \Ref{V.4}. Let us prove (iii)$\To$(iv). Indeed,
we know that $\varprojlim_i Q_i$ is nonempty as a projective
limit of nonempty finite sets. Let $a$ be a point of
$\varprojlim_i Q_i$; it defines a homomorphism of spaces with
operators
$$
\pi\to Q
$$
which is \emph{surjective}, since for every $i$ the composite $\pi\to
Q\to Q_i$ is, because $\pi$ is transitive on $Q_i$ by virtue
of~\Ref{V.5.3}. If therefore $H$ is the stabilizer subgroup of
$a$, one obtains an isomorphism $\pi/H\isomto Q$. The implications
(iv)$\To$(v) and (v)$\To$(i) are again
trivial.

\begin{proposition}
\label{V.5.6}
Let $\cal{C}$ be a Galois category, $P$ a fundamental
pro-object of~$\cal{C}$ and $F$ the associated fundamental
functor. Let $P'=(P'_i)_{i\in I}$ be a pro-object of~$\cal{C}$, put
in normal form, and $F'$ the associated functor $F'(X)=\Hom(P',X)$
from~$\cal{C}$ into~$\Ens$. Equivalent conditions:
\begin{enumerate}
\item[(i)] $P'\simeq P$, or again $F'\simeq F$.
\item[(ii)] $P'$ is fundamental, or again $F'$ is
fundamental.
\item[(iii)] $F'$ transforms a sum of two objects into a sum, and
$X\not\simeq\emptyset_{\cal{C}}$ implies $F(X)\ne\emptyset$.
\item[(iv)]
% original p. 130
The objects of~$\cal{C}$ that are connected and $\ne\emptyset_{\cal{C}}$ are
exactly the objects isomorphic to a~$P'_i$.
\end{enumerate}
\end{proposition}

We have trivially (i)$\To$(iii) and (i)$\To$(ii),
moreover (ii)$\To$(iv) by virtue of~\Ref{V.5.4} (applied
to $P'$ instead of~$P$). Moreover (iii) or (iv) implies by virtue
of~\Ref{V.5.5} that $P'$ is of the form $\pi/H$ where $H$ is a
closed subgroup of~$\pi$. In case (iii), there exists for every
open invariant subgroup $\pi'$ of~$\pi$ a $\pi$-homomorphism
$P'=\pi/H\to \pi/\pi'$, whence $H\subset \pi'$, whence $H=(0)$ and
consequently (i), qed.

\begin{corollary}
\label{V.5.7}
Let $\cal{C}$ be a Galois category. The fundamental
pro-objects are isomorphic, the fundamental functors are
isomorphic.
\end{corollary}

In other words, \emph{the category of fundamental functors
is a connected groupoid~$\Gamma$},
which one may call the \emph{fundamental groupoid}
of the Galois category~$\cal{C}$. If $\cal{C}=\cal{C}(\pi)$,
the group of automorphisms of an object of the fundamental groupoid
is isomorphic to $\pi$, this isomorphism being well
determined up to inner automorphism. (N.B. one
calls a \emph{groupoid} a category in which all the
morphisms are isomorphisms, a \emph{connected} groupoid a
groupoid all of whose objects are isomorphic). The fundamental
pro-objects of~$\cal{C}$ form a connected groupoid
equivalent to the \emph{opposite} of the fundamental
groupoid. If $F, F'$ are two fundamental functors, associated
with fundamental pro-objects $P, P'$, then $\Hom(F, F')=\Isom(F,
F')$ is sometimes denoted $\pi_{F', F}$ and plays the role of a
``set of \emph{classes of paths} from~$F$ to~$F'$''. In
particular $\pi_{F, F}=\pi_F$ is none other than the \emph{fundamental
group of~$\cal{C}$} at $F$ constructed in the preceding
no. As for the pro-object $P$ associated with~$F$, it
plays the role of a \emph{universal covering at~$F$} of
the final object $e_{\cal{C}}$ of~$\cal{C}$.

It may be convenient to have a description of~$\cal{C}$ (up to
equivalence) in terms of its fundamental groupoid
$\Gamma$, without going through a choice of a particular object $F$ of the
latter. Now to every object $X$ of~$\cal{C}$ is associated the
functor $E_X$ on the fundamental groupoid, defined by
$$
E_X(F)=F(X),
$$
with values in~$\Ens$. (Such a functor is known in topology
under the name of ``local system'' on the groupoid)
$F(X)=E_X(F)$ may be called the \emph{fiber} of~$X$ at $F$,
and the functor $E_X$ the fiber-functor associated with~$X$. The
functor $E_X$ has the following
% original p. 131
property:
\emph{for every $F$, $E_X(F)$ is a finite set
on which the topological group $\pi_F=\Aut(F)$ operates continuously}.
For a
given covariant functor $\xi$ from the fundamental groupoid
into~$\Ens$, the preceding condition is moreover equivalent
to the same condition for \emph{one} arbitrary fixed $F$.
This being so:

\begin{proposition}
\label{V.5.8}
The functor $X\mto E_X$ is an equivalence of the category
$\cal{C}$ with the category of covariant functors from the
fundamental groupoid $\Gamma$ of~$\cal{C}$ into~$\Ens$, which
satisfy the condition
highlighted above.
\end{proposition}

Indeed, let $F_0$ be an object of the fundamental groupoid, and let
$\pi_0=\Aut(F_0)$; then the functor $\xi\mto \xi(F_0)$ is an
equivalence of the second category considered
in~\Ref{V.5.8} with the category $\cal{C}(\pi_0)$, as one
sees at once. On the other hand, the composite of the latter and
of the functor $X\mto E_X$ is the natural equivalence
$\cal{C}\to\cal{C}(\pi_0)$. It follows that the functor
$X\mto E_X$ itself is an equivalence.

\begin{corollary}
\label{V.5.9}
The category $\Prodash\cal{C}$
is canonically equivalent to the category of covariant
functors $\xi$ from the fundamental groupoid $\Gamma$ into the
category of topological spaces, satisfying the condition:
for every object $F$ of~$\Gamma$, $\xi(F)$ is a compact totally
disconnected space with topological group $\pi_F$
of operators.
\end{corollary}

Here again, one can verify this condition on~$\xi$; it suffices
to verify it for \emph{one}~$F$. The proof is the
same as for~\Ref{V.5.8}.

\begin{remark}
\label{V.5.10}
Let $(F_s)_{s\in S}$ be a family
of objects
of the fundamental groupoid~$\Gamma$. Let us set, for $s, s'\in S$:
$$
\Hom(s, s')=\Hom(F_s, F_{s'})
$$
so that $S$ itself becomes a connected groupoid and
the map $s\mto F_s$ a fully faithful functor from~$S$
into $\Gamma$, call it~$f$. Considering then the functor $X\mto
E_X\circ f$ from~$\cal{C}$ into the functor category
$\Hom(S,\Ens)$, one obtains a variant of~\Ref{V.5.8}
(and~\Ref{V.5.9}) with $\Gamma$ replaced by~$S$. The statement
thus obtained reduces to Theorem~\Ref{V.4.1} when $S$
is reduced to a point, and is none other than~\Ref{V.5.8}
itself if $S$ is the set of objects of~$\Gamma$.
\end{remark}

We are going to use~\Ref{V.5.9} to define a canonical
pro-object of~$\cal{C}$. For this, we consider the functor
from~$\Gamma$ into the category of topological spaces
% original p. 132
(and even of topological groups):
$$
f\colon F\mto \Aut(F)=\pi_F.
$$
This functor satisfies the condition considered in~\Ref{V.5.8};
the space with operators $f(F)$ under $\pi_F$ is none other than
$\pi_F$, considered as a space with operators under
itself by inner automorphisms. Hence the functor
$f$ corresponds to a pro-object of~$\cal{C}$ determined up to
unique isomorphism, which is even a pro-group of
$\cal{C}$ and which is called the \emph{fundamental pro-group
of~${\cal C}$},
playing the role of a local system of fundamental
groups. It is therefore a pro-group $\Pi$ of~$\cal{C}$ defined
by the condition that one have an isomorphism functorial in $F$
$$
F(\Pi)\simeq\pi_F.
$$

If $X$ is an arbitrary pro-object of~$\cal{C}$, one has a canonical
morphism
$$
\Pi\times X \to X
$$
which makes~$X$ an object with a group of operators on the left
$G$ in $\Prodash\cal{C}$. For this it suffices to note that for $F$
variable, one has a canonical map
$$
\Pi(F)\times X(F) \to X(F)
$$
\ie $$
\Aut(F)\times E_X(F) \to E_X(F), \quad \text{or} \quad \pi_F\times
F(X)\to F(X)
$$
which is functorial in~$F$. It is also functorial in $X$, hence
for every morphism $X\to Y$ of pro-objects, the corresponding
diagram
$$
\xymatrix{
\Pi \times X \ar[r]\ar[d]& X\ar[d] \\ \Pi \times Y \ar[r]& Y
}
$$
is commutative.

\begin{remark}
\label{V.5.11}
One must be careful not to confuse a fundamental pro-object $P$ (which is not
endowed with a group structure, and is connected) with the fundamental
pro-group (which is a pro-\emph{group}, and in general not
connected). Precisely, $G$ is connected if
and
only if $\pi_F$ operating on itself by inner
automorphisms is transitive, \ie if $\pi$ is reduced to
the neutral element, or again $\cal{C}$ equivalent to the
category of finite sets. Another
% original p. 133
essential difference is that $G$ is determined up to
unique isomorphism, and $P$ is determined only up to
non-unique isomorphism.

Let $E$ be a finite set and let us consider the constant functor on
the groupoid~$\Gamma$ with value $E$: by virtue
of~\Ref{V.5.8} it defines an object of~$\cal{C}$, denoted $E_{\cal{C}}$, which
may also be interpreted as the sum of~$E$ copies of
the final object $e_{\cal{C}}$ of~$\cal{C}$. One may consider
$E_{\cal{C}}$ as a functor in $E$, from the category of
finite sets into the category $\cal{C}$, and this functor is
\emph{exact}, hence transforms finite groups into $\cal{C}$-groups,
etc. If therefore $X$ is an object of~$\cal{C}$ on which the finite group
$G$ operates on the right, we see that one may consider $X$
as an object of~$\cal{C}$ having a $\cal{C}$-group of operators
on the right~$G_{\cal{C}}$. We shall therefore say, by extension of a
general terminology relative to objects with
$\cal{C}$-groups of operators, that $X$ is \emph{formally
principal homogeneous}
under $G$ if $X$ is formally principal homogeneous under
$G_{\cal{C}}$, \ie if the canonical morphism
$$
X\times G_{\cal{C}} \to X\times X
$$
deduced from the operation of~$G_{\cal{C}}$ on~$X$ on the right,
is an isomorphism. We say that $X$ is \emph{principal homogeneous}
under $G$ if it is so under $G_{\cal{C}}$, \ie if it is so formally,
and if moreover the quotient $X/G=X/G_{\cal{C}}$ is~$e_{\cal{C}}$. If one
fixes a fundamental functor, whence an equivalence of
$\cal{C}$ with a category $\cal{C}(\pi)$, $X$ corresponds to
a set on which $\pi$ operates on the left continuously,
say $E=F(X)$. Making $G$ operate on~$X$ on the right then amounts
to making $G$ operate on the set $E$ on the right, in such a way that the operations of~$G$ commute with those of~$\pi$. One
then sees at once that $X$ is principal homogeneous under
$G$ if and only if the set $E$ is a principal
homogeneous space
under $G$,
\ie if and only if $G$ operates on it in a simply transitive way. (Moreover $X$ is formally
principal homogeneous if and only if
$E$ is principal homogeneous \emph{or} empty). Comparing
with~\Ref{V.5.3}, we see that if $X$ is principal homogeneous under
$G$ \emph{and connected}, then the given homomorphism from~$G$ into the
group opposite to $\Aut(X)$ is an \emph{isomorphism}; and
moreover, for an object $X$
of~$\cal{C}$
to be connected and principal homogeneous under the group opposite to
$\Aut(X)$, it is necessary and sufficient, with the notation of \No \Ref{V.4}, that it
be isomorphic to a \emph{Galois} $P_i$. In the typical case
$\cal{C}=\cal{C}(\pi)$, this means that $X$ is isomorphic to a
quotient of~$\pi$ by an invariant subgroup.

Let us still suppose given a fundamental functor~$F$. Then the
datum of an
% original p. 134
$X$ principal homogeneous under a finite group $G$ operating on
the right, and of a point $a\in F(X)$, is equivalent to the
datum of a homomorphism from~$\pi$ into the group~$G$. Indeed,
to such a homomorphism one makes correspond the set $E=G$,
making~$\pi$ operate on it on the left by means of the given homomorphism
$\pi\to G$ and the left translations of~$G$, and
making $G$ operate on it on the right by right translation, the
marked point $a$ of~$E$ being the unit element
of~$G$. Thanks to the foregoing, one thus obtains in an
essentially unique way every triple $(X, G, a)$ having the
properties considered above,
since a
principal homogeneous set with marked point under a group $G$ is identified
with the latter. In this way, one has a direct geometric
interpretation of the functor $G\mto \Hom(\pi, G)$ from the
category of finite groups into~$\Ens$, a functor which is
pro-representable by means of~$\pi$, and whose
consideration would therefore furnish another construction of the group
$\pi$ associated with~$F$.
\end{remark}
